\section{Thiết kế thực nghiệm}
\label{sec:expdesign}

\subsection{Mục tiêu và câu hỏi nghiên cứu}

Phần thực nghiệm được thiết kế để trả lời ba câu hỏi nghiên cứu, đồng thời kiểm
chứng tính hiệu quả của pipeline đã đề xuất:

\begin{description}[leftmargin=2.4em,itemsep=3pt]
  \item[\textbf{CH1.}] ResNet18 có vượt được CNN baseline hay không, và nếu có
        thì độ sâu cùng kết nối phần dư đóng góp ra sao?
  \item[\textbf{CH2.}] Augmentation tác động thế nào tới hiệu năng -- không chỉ ở
        mức tổng thể mà còn ở mức từng lớp?
  \item[\textbf{CH3.}] Ensemble hai mô hình ResNet có khai thác được tính bù trừ
        lỗi giữa chúng để cho kết quả tốt hơn mọi mô hình đơn lẻ hay không?
\end{description}

\subsection{Môi trường thực hiện}

Toàn bộ thực nghiệm được tiến hành trong học kỳ 252 theo yêu cầu môn học. Mã
nguồn được tổ chức thành các file Python trong thư mục \texttt{src/}, chạy được
trên cả Google Colab (GPU) và máy cá nhân. Việc tách riêng phần tiền xử lý,
\texttt{Dataset}, mô hình, huấn luyện và đánh giá giúp project dễ kiểm thử và mở
rộng. Mọi lần chạy đều được ghi log bằng Weights \& Biases.

\subsection{Chuẩn đầu vào cố định}

Để khác biệt về kết quả chỉ đến từ mô hình chứ không từ dữ liệu, các yếu tố sau
được giữ nguyên trong mọi cấu hình:

\begin{itemize}[leftmargin=1.4em,itemsep=1pt]
  \item bộ dữ liệu UrbanSound8K, chia fold cố định (test = fold 1, val = fold 2);
  \item đầu vào là log-Mel spectrogram $128 \times 173$;
  \item tần số 22{,}05\,kHz, độ dài 4{,}0 giây ($88200$ mẫu);
  \item tham số đặc trưng $n_{\text{fft}}=1024$, $\text{hop}=512$,
        $n_{\text{mels}}=128$, $\text{top\_db}=80$;
  \item bộ chỉ số đánh giá gồm Accuracy, Macro-F1 và confusion matrix.
\end{itemize}

\subsection{Bốn cấu hình thực nghiệm}

Đề tài triển khai và đánh giá bốn cấu hình. Ba cấu hình đầu là các mô hình được
huấn luyện độc lập; cấu hình thứ tư không huấn luyện thêm mà kết hợp hai mô hình
ResNet đã có.

\begin{description}[leftmargin=2.0em,itemsep=3pt]
  \item[\textbf{C1 -- CNN baseline.}] Mô hình cơ sở, làm mốc tham chiếu.
  \item[\textbf{C2 -- ResNet18.}] ResNet18 chỉnh cho đầu vào spectrogram một
        kênh, không pretrained, không augmentation.
  \item[\textbf{C3 -- ResNet18 + augmentation.}] Cùng backbone ResNet18 nhưng bổ
        sung SpecAugment và MixUp ở nhánh train, dung lượng mô hình lớn hơn.
  \item[\textbf{C4 -- Ensemble.}] Trung bình xác suất của C2 và C3 (đều kèm TTA).
\end{description}

Bảng~\ref{tab:config} tóm tắt siêu tham số huấn luyện của ba mô hình. Các giá trị
này được chọn dựa trên thử nghiệm sơ bộ và đặc điểm của từng cấu hình -- ví dụ,
C3 dùng dropout và weight decay lớn hơn vì augmentation đã đưa thêm nhiễu, nên mô
hình cần dung lượng lớn hơn nhưng cũng cần được điều hòa mạnh hơn.

\begin{table}[H]
\centering
\caption{Siêu tham số huấn luyện của ba mô hình.}
\label{tab:config}
\small
\begin{tabular}{lccc}
\toprule
\textbf{Siêu tham số} & \textbf{C1 -- CNN} & \textbf{C2 -- ResNet18} & \textbf{C3 -- ResNet18+aug}\\
\midrule
Số epoch tối đa      & 30                & 40                 & 50\\
Optimizer            & Adam              & AdamW              & AdamW\\
Learning rate        & $1\!\times\!10^{-3}$ & $5\!\times\!10^{-4}$ & $3\!\times\!10^{-4}$\\
Weight decay         & $1\!\times\!10^{-4}$ & $1\!\times\!10^{-4}$ & $1\!\times\!10^{-3}$\\
Lập lịch LR          & ReduceLROnPlateau & Cosine             & Cosine\\
Label smoothing      & 0{,}10            & 0{,}05             & 0{,}05\\
Dropout              & 0{,}30            & 0{,}30             & 0{,}50\\
Base channels        & --                & 48                 & 64\\
Early stopping       & 5                 & 8                  & 10\\
Augmentation         & Không             & Không              & SpecAugment + MixUp\\
\bottomrule
\end{tabular}
\end{table}

\subsection{Quy trình huấn luyện và đánh giá}

Mỗi mô hình được huấn luyện trên tập train và theo dõi trên tập validation sau
từng epoch. Checkpoint có Accuracy validation cao nhất được lưu lại; chỉ checkpoint
này -- chứ không phải mô hình ở epoch cuối -- mới được dùng để đánh giá test. Quy
trình này đảm bảo kết quả báo cáo không phải là sản phẩm của một epoch may mắn,
cũng không phải kết quả của việc tinh chỉnh trực tiếp trên tập test.

Khi đánh giá, mô hình được chạy trên toàn bộ tập test (fold 1). Với C2, C3 và C4,
test-time augmentation được áp dụng như mô tả ở Mục~\ref{sec:method}. Đầu ra gồm
Accuracy, Macro-F1, classification report theo từng lớp và confusion matrix --
tất cả được lưu lại để phục vụ phân tích.

\subsection{Tiêu chí so sánh}

Việc so sánh dựa trên cả khía cạnh định lượng lẫn định tính. Về định lượng,
Accuracy cho biết mô hình nào dự đoán đúng nhiều nhất trên toàn tập test, còn
Macro-F1 cho biết mô hình nào cân bằng nhất giữa các lớp. Hai chỉ số này không
phải lúc nào cũng đồng thuận, và chính sự khác biệt giữa chúng là một phần thông
tin đáng giá. Về định tính, confusion matrix và classification report giúp truy
ngược: vì sao một mô hình có Accuracy cao nhưng Macro-F1 thấp, hay augmentation
thực sự đã thay đổi điều gì ở mức từng lớp.

\subsection{Tính tái lập}

Đề tài chú trọng tính tái lập ở ba mức. Mã nguồn được chia thành các file chức
năng rõ ràng. Pipeline tiền xử lý được cố định và mô tả minh bạch, kèm cơ chế
feature cache để mọi lần chạy dùng đúng cùng một đặc trưng. Toàn bộ kết quả thí
nghiệm được log bằng W\&B, kèm seed cố định ở mức hợp lý. Nhờ vậy, người đọc có
thể hiểu kết quả được tạo ra như thế nào và chạy lại khi cần.
